% Demo/reference slide: every color defined by whichever per-university
% preamble is currently \input{}ed (Slidemaster_preamble_uni-graz.tex,
% Slidemaster_preamble_mci.tex, or Slidemaster_preamble_uni-ibk.tex) -- this
% slide renders the same way under all three, showing each one's own actual
% palette, since the color *names* (c_primary_dark, c_secondary_blue, ...)
% are shared machinery, only the underlying hex values differ per
% institution. Labels below are the full color name (e.g. the first swatch
% is \texttt{c\_primary\_dark}), not a "suffix after a common prefix" the
% way this slide worked when there was only ever one institution -- the new
% names don't share one single prefix across all four tiers
% (c_primary_*/c_secondary_*/c_tertiary_*/c_custom_*), so the swatch list
% below spells out each full name directly instead of looping a suffix onto
% one fixed prefix. Sorted by relevance: the four primary colors first
% (dark, highlight, gray, white), then the secondary colors (six hues, each
% paired with its tertiary/lighter variant), and finally the three custom
% colors last -- not part of any institution's official palette, just
% supporting colors this project defines itself (see each color's own
% \definecolor comment in whichever Slidemaster_preamble_*.tex is current).
%
% Copy from here when starting a new slide; see README.md for the
% naming/versioning convention.
%% Included in decks: test-slides
%%

\setSlideVersion{1-0-0}
\begin{standardframe}{Slidemaster color palette}
	\footnotesize
	Every color defined by the current institution's preamble. Labels below
	are the color's full macro name -- e.g. the first swatch is
	\texttt{c\_primary\_dark}.

	\vspace{0.8em}
	\centering
	% Each color name contains literal underscores (e.g. "c_primary_dark")
	% -- outside math mode, a bare "_" is a special character that LaTeX
	% reads as "Missing $ inserted", auto-recovering by subscripting just
	% the next token, which prints as "c_p rimary_d ark" rather than the
	% literal name. Fixed by driving the loop off a name/label *pair* per
	% entry (pgf's \foreach \a/\b in {...} two-variable syntax) rather than
	% trying to print the raw name directly: \name (unescaped) still goes
	% into fill=\name for the color lookup, but \dispname carries a
	% separately-written, already-\_-escaped copy of the same string for
	% display -- sidesteps the underscore-catcode question entirely, no
	% runtime \catcode trick needed. Named \dispname rather than the more
	% obvious \label: a first attempt using \label rendered as empty text
	% throughout, silently -- \label collides with plain LaTeX's own
	% cross-referencing \label{} command, and \foreach's local redefinition
	% of it didn't behave the way a fresh, previously-undefined loop
	% variable name does.
	\begin{tikzillustration}[
			swatch/.style={rectangle, draw=c_primary_gray, minimum width=1.75cm, minimum height=0.85cm, inner sep=0pt},
			swlabel/.style={font=\fontsize{4.3}{5}\selectfont\ttfamily, align=center, text width=1.95cm}
		]{zzm-tikz}
		\foreach \name/\dispname [count=\i] in {%
			c_primary_dark/c\_primary\_dark, c_primary_highlight/c\_primary\_highlight, c_primary_gray/c\_primary\_gray, c_primary_white/c\_primary\_white,%
			c_secondary_violet/c\_secondary\_violet, c_tertiary_violet/c\_tertiary\_violet, c_secondary_blue/c\_secondary\_blue, c_tertiary_blue/c\_tertiary\_blue, c_secondary_green/c\_secondary\_green, c_tertiary_green/c\_tertiary\_green,%
			c_secondary_red/c\_secondary\_red, c_tertiary_red/c\_tertiary\_red, c_secondary_orange/c\_secondary\_orange, c_tertiary_orange/c\_tertiary\_orange, c_secondary_yellow/c\_secondary\_yellow, c_tertiary_yellow/c\_tertiary\_yellow,%
			c_tertiary_gray/c\_tertiary\_gray, c_custom_gray/c\_custom\_gray, c_custom_red/c\_custom\_red, c_custom_green/c\_custom\_green, c_custom_brown/c\_custom\_brown%
		} {
			\pgfmathsetmacro{\col}{mod(\i-1,7)}
			\pgfmathsetmacro{\row}{floor((\i-1)/7)}
			\coordinate (sw) at (\col*2.0cm, -\row*1.45cm);
			\node[swatch, fill=\name] at (sw) {};
			\node[swlabel, anchor=north] at ($(sw)+(0,-0.47cm)$) {\dispname};
		}
	\end{tikzillustration}
\end{standardframe}
